%==============================================================
%  COMP34812 NLU Coursework Poster
%  Track A: Natural Language Inference (NLI)
%  Format: A1 Landscape (compliant with coursework spec)
%
%  Layout:
%    [Title bar — full width]
%    [Task & Dataset & Project Spec — full width]
%    [Solution B: RWKV-7  |  Solution C: Transformer Ensemble]
%    [Comparison + Limitations & Ethics — full width]
%==============================================================
\documentclass[final,t]{beamer}
\usepackage[orientation=landscape,size=a1,scale=1.15]{beamerposter}
\usepackage[utf8]{inputenc}
\usepackage[T1]{fontenc}
\usepackage{lmodern}
\usepackage{booktabs}
\usepackage{array}
\usepackage{amsmath,amssymb}
\usepackage{graphicx}
\usepackage[table]{xcolor}
\usepackage{tikz}
\usetikzlibrary{positioning,arrows.meta}
\usepackage{fontawesome5}

%----- Colour scheme ------------------------------------------
\definecolor{uomPurple}{RGB}{87,36,127}
\definecolor{uomYellow}{RGB}{233,183,23}
\definecolor{softGrey}{RGB}{245,245,248}
\definecolor{accentRed}{RGB}{190,40,60}
\definecolor{accentGreen}{RGB}{30,130,76}
\definecolor{rwkvBlue}{RGB}{28,91,163}
\definecolor{ensembleOrange}{RGB}{198,92,25}

%----- Beamer theme tweaks ------------------------------------
\usetheme{default}
\setbeamercolor{block title}{fg=white,bg=uomPurple}
\setbeamercolor{block body}{fg=black,bg=softGrey}
\setbeamerfont{block title}{size=\large,series=\bfseries}
\setbeamertemplate{itemize item}{\color{uomPurple}$\blacktriangleright$}
\setbeamertemplate{itemize subitem}{\color{uomYellow}$\bullet$}
\setbeamertemplate{navigation symbols}{}

% Helper: coloured block titles for the two solution columns
\newenvironment{rwkvblock}[1]{%
  \setbeamercolor{block title}{fg=white,bg=rwkvBlue}%
  \begin{block}{#1}%
}{\end{block}}
\newenvironment{ensembleblock}[1]{%
  \setbeamercolor{block title}{fg=white,bg=ensembleOrange}%
  \begin{block}{#1}%
}{\end{block}}

\title{}
\author{}
\date{}

\begin{document}
\begin{frame}[t]
\vspace*{-0.4cm}

%==============================================================
%  TITLE BAR  (full width — matches the body layout)
%==============================================================
\begin{columns}[t,totalwidth=\paperwidth]
\begin{column}{0.97\paperwidth}
\begin{tikzpicture}
    \node[fill=uomPurple,text=white,minimum width=0.97\paperwidth,
          minimum height=6.2cm,inner sep=0.55cm] (title) {
          \begin{minipage}{0.95\paperwidth}
              \centering
              {\Huge\bfseries Natural Language Inference:
               Recurrent LLM vs.\ Transformer Ensemble}\\[0.45cm]
              {\LARGE Serhii Tupikin \;\textbar\; Ilya Maltsev
               \;\textbar\; Julius Vander Arend
               \;\;$\diamond$\;\; Group CW2
               \;\;$\diamond$\;\; Track A: NLI}\\[0.15cm]
              {\Large COMP34812 Natural Language Understanding \;\textbar\;
               The University of Manchester, 2025--26}
          \end{minipage}
    };
    \draw[uomYellow,line width=6pt]
         ([yshift=-4mm]title.south west) --
         ([yshift=-4mm]title.south east);
\end{tikzpicture}
\end{column}
\end{columns}
\vspace*{-1.2cm}

%==============================================================
%  FULL-WIDTH : Task, Dataset & Project Spec
%==============================================================
\begin{columns}[t,totalwidth=\paperwidth]
\begin{column}{0.97\paperwidth}

\begin{block}{\faBullseye\ \ Task \& Dataset}
\vspace{0.4em}
\begin{columns}[c]

%-- Sub-col 1: Task definition --
\begin{column}{0.48\textwidth}
\centering
\textbf{Natural Language Inference (NLI).}
Given a \emph{premise} $p$ and a \emph{hypothesis} $h$,
decide whether $h$ is entailed by $p$.\\[0.2em]
\begin{tabular}{@{}r@{\;}l@{}}
\textbf{Input:}  & sentence pair $(p,h)$ \\
\textbf{Output:} & \emph{entailment} or \emph{contradiction}
                   — binary label $\{1,\,0\}$ \\
\textbf{Track:}  & A — pairwise sequence classification \\
\end{tabular}
\end{column}

%-- Sub-col 2: Dataset stats --
\begin{column}{0.48\textwidth}
\centering
\textbf{NLU Shared-Task NLI corpus}\\
\emph{no external datasets allowed}\\[0.2em]
\renewcommand{\arraystretch}{1.05}
\begin{tabular}{lr}
\toprule
\textbf{Split} & \textbf{Pairs} \\
\midrule
Training        & 24{,}432 \\
Development     & 6{,}736  \\
Test (hidden)   & 3{,}302  \\
\bottomrule
\end{tabular}
\end{column}

\end{columns}
\vspace{-0.3em}
\end{block}

\end{column}
\end{columns}

\vspace{0.15cm}

%==============================================================
%  TWO-COLUMN BODY : RWKV  |  Transformer Ensemble
%  Raw side-by-side minipages with a hard fixed height — this
%  GUARANTEES both columns end at the same Y so the bottom row
%  is perfectly aligned.  No \begin{columns} quirks.
%==============================================================
\newlength{\solncolheight}
\setlength{\solncolheight}{31cm}
\newlength{\solncolwidth}
\setlength{\solncolwidth}{0.485\paperwidth}

\noindent
%--------------------------------------------------------------
%  LEFT COLUMN : RWKV-7 "Goose"  (Category B)
%--------------------------------------------------------------
\begin{minipage}[t][\solncolheight][t]{\solncolwidth}

\vspace{1.0em}
\begin{center}
\colorbox{rwkvBlue}{\color{white}\large\bfseries
\hspace{0.4em}\faCrow\ \ Solution B: RWKV-7 ``Goose''\hspace{0.4em}}
\end{center}

\begin{rwkvblock}{Model Overview}
\textbf{An attention-free, recurrent 7.2B-parameter LLM.}
\begin{itemize}
    \item \textbf{Architecture:} RWKV-7 — an advanced RNN
          reaching transformer-level quality without
          quadratic self-attention.
    \item \textbf{Core innovation:} the \emph{Delta Rule}
          replaces softmax attention with a recurrent state
          update, removing the $O(n^2)$ bottleneck.
    \item \textbf{Complexity:}
          \;\faClock\ \textbf{Time:} $\mathcal{O}(n)$ \quad
          \;\faMemory\ \textbf{Space:} $\mathcal{O}(1)$
    \item \textbf{Why it fits NLI:} long-range, streaming-friendly
          reasoning over $(p,h)$ pairs at a fraction of transformer
          inference cost.
\end{itemize}
\end{rwkvblock}

\begin{rwkvblock}{Training Strategy: LoRA Fine-Tuning}
\textbf{Low-Rank Adaptation} tailored for binary NLI classification.
\begin{center}
\renewcommand{\arraystretch}{1.15}
\begin{tabular}{ll}
\toprule
\textbf{Rank / Alpha}     & $r = 32$ \;\;/\;\; $\alpha = 64$ \\
\textbf{Trainable params} & $\sim$83.8M \;(\textbf{1.15\%} of total) \\
\textbf{Frozen params}    & $\sim$7.12B \\
\textbf{Hardware}         & 1$\times$ RTX 5090 \\
\textbf{Training time}    & \textbf{$\sim$2 h} for 4 epochs \\
\bottomrule
\end{tabular}
\end{center}
\textbf{Why LoRA?} Small low-rank matrices are added into the
frozen RWKV-7 model — updating only $\sim$1.15\% of its weights
— to fine-tune it for NLI while keeping its general language
abilities intact.
\end{rwkvblock}

\begin{rwkvblock}{Results on Dev Set}
\centering
\renewcommand{\arraystretch}{1.25}
\begin{tabular}{lccc}
\toprule
 & \textbf{F1} & \textbf{Accuracy} & \textbf{MCC} \\
\midrule
Baseline (Cat.\ B) & \textit{[\,--\,]} & \textit{[\,--\,]} & \textit{[\,--\,]} \\
\rowcolor{softGrey}
\textcolor{accentGreen}{\textbf{\faTrophy\ \ Our RWKV-7 (LoRA)}} &
\textcolor{accentGreen}{\textbf{0.9419}} &
\textcolor{accentGreen}{\textbf{0.9399}} &
\textcolor{accentGreen}{\textbf{0.8796}} \\
\bottomrule
\end{tabular}
\end{rwkvblock}

\begin{rwkvblock}{Key Takeaways}
\begin{itemize}
    \item \textbf{Delta rule is all you need}: the attention-free
          recurrent backbone matches transformer-level quality at
          \emph{linear-time} inference cost.
    \item \textbf{Matches a heavy transformer ensemble}: despite
          being recurrent, RWKV-7 reached the same task-level
          performance as our compute-heavy Transformer Ensemble
          (Solution C) on this NLI task.
    \item \textbf{Scale vs.\ efficiency}: the lighter 2.9B
          variant came \emph{close but not as good} even after 7
          epochs (vs.\ 4 for 7.2B) — we chose 7.2B because the
          goal was maximum achievable performance, not minimum
          compute.
\end{itemize}
\end{rwkvblock}

\vfill
\end{minipage}%
\hfill
%--------------------------------------------------------------
%  RIGHT COLUMN : Transformer Ensemble  (Category C)
%--------------------------------------------------------------
\begin{minipage}[t][\solncolheight][t]{\solncolwidth}

\vspace{1.0em}
\begin{center}
\colorbox{ensembleOrange}{\color{white}\large\bfseries
\hspace{0.4em}\faLayerGroup\ \ Solution C: Transformer Ensemble\hspace{0.4em}}
\end{center}

\begin{ensembleblock}{Model Overview}
\textit{[TODO: confirm exact base models and combination strategy]}
\begin{itemize}
    \item \textbf{Architecture:} ensemble of
          \textit{[e.g.\ DeBERTa-v3-large, RoBERTa-large, BERT-large]}
          fine-tuned independently on the NLI corpus.
    \item \textbf{Combination:}
          \textit{[e.g.\ soft-voting / logit averaging / stacking]}
          over the per-model probability outputs.
    \item \textbf{Input format:}
          \texttt{[CLS] premise [SEP] hypothesis [SEP]}.
    \item \textbf{Why an ensemble:} reduces variance of individual
          fine-tunes, captures diverse lexical/syntactic priors
          from different pre-training corpora.
\end{itemize}
\end{ensembleblock}

\begin{ensembleblock}{Training Strategy}
\textit{[TODO: fill in once you have the final config]}
\begin{center}
\renewcommand{\arraystretch}{1.15}
\begin{tabular}{ll}
\toprule
\textbf{Base models}     & \textit{[model\_1, model\_2, model\_3]} \\
\textbf{Fine-tune method}& Full fine-tune / head-only / adapters \\
\textbf{Optimiser}       & AdamW, lr $= \textit{[\dots]}$ \\
\textbf{Batch size}      & \textit{[\dots]} \\
\textbf{Epochs}          & \textit{[\dots]} per model \\
\textbf{Hardware}        & \textit{[e.g. RTX 4090 / A100]} \\
\textbf{Total train time}& \textit{[\dots]} \\
\bottomrule
\end{tabular}
\end{center}
Each model is validated independently; the final ensemble
weights are selected on the development set to optimise F1.
\end{ensembleblock}

\begin{ensembleblock}{Results on Dev Set}
\noindent
% --- LEFT: results table ---------------------------------------
\begin{minipage}[t]{0.47\linewidth}
\centering
\vspace{0pt}%
\renewcommand{\arraystretch}{1.25}
\begin{tabular}{lccc}
\toprule
 & \textbf{F1} & \textbf{Acc.} & \textbf{MCC} \\
\midrule
Baseline (C)  & \textit{[-]} & \textit{[-]} & \textit{[-]} \\
Single best   & \textit{[-]} & \textit{[-]} & \textit{[-]} \\
\rowcolor{softGrey}
\textbf{Our Ensemble} & \textit{[-]} & \textit{[-]} & \textit{[-]} \\
\bottomrule
\end{tabular}

\vspace{0.4em}
\small\textit{[TODO: add headline finding once
numbers are in.]}
\end{minipage}\hfill
% --- RIGHT: chart placeholder ----------------------------------
\begin{minipage}[t]{0.50\linewidth}
\centering
\vspace{0pt}%
% Drop your chart here:
%   \includegraphics[width=\linewidth]{figures/ensemble_f1.pdf}
\fbox{\begin{minipage}[c][3.8cm][c]{\dimexpr\linewidth-2\fboxsep\relax}
\centering
\small\textit{[TODO: chart — e.g.\ per-model F1 vs.\ ensemble,
or confusion matrix. Replace this box with
\texttt{\textbackslash includegraphics\{\dots\}}.]}
\end{minipage}}
\end{minipage}
\end{ensembleblock}

\begin{ensembleblock}{Key Takeaways}
\begin{itemize}
    \item \textit{[TODO: ensemble adds X F1 points over the single
          best transformer.]}
    \item \textit{[TODO: diverse base models contribute the most
          to the gain (low prediction correlation).]}
    \item \textit{[TODO: diminishing returns beyond N=? members.]}
\end{itemize}
\end{ensembleblock}

\vfill
\end{minipage}%

\vspace{0.3cm}

%==============================================================
%  BOTTOM ROW : Comparison + Error Analysis  |  Limitations
%  Same raw-minipage pattern for guaranteed alignment.
%==============================================================
\newlength{\bottomrowheight}
\setlength{\bottomrowheight}{11cm}

\noindent
\begin{minipage}[t][\bottomrowheight][t]{\solncolwidth}

\begin{block}{\faBalanceScaleLeft\ \ Head-to-Head Comparison (Dev)}
\centering
\renewcommand{\arraystretch}{1.25}
\begin{tabular}{lcccc}
\toprule
\textbf{System} & \textbf{Params (train)} & \textbf{F1}
                & \textbf{Acc.} & \textbf{MCC} \\
\midrule
\textcolor{rwkvBlue}{\textbf{RWKV-7 + LoRA}} &
$\sim$83.8M & \textbf{0.9419} & 0.9399 & 0.8796 \\
\textcolor{ensembleOrange}{\textbf{Transformer Ensemble}} &
\textit{[\,--\,]} & \textit{[\,--\,]} & \textit{[\,--\,]}
                  & \textit{[\,--\,]} \\
\bottomrule
\end{tabular}

\vspace{0.3em}
\small A recurrent LLM with 1.15\% tunable parameters competes
directly with a stacked transformer ensemble — the core
headline of this project.
\end{block}

\vfill
\end{minipage}%
\hfill
\begin{minipage}[t][\bottomrowheight][t]{\solncolwidth}

\begin{block}{\faBalanceScale\ \ Limitations \& Ethical Considerations}
\begin{itemize}
    \item \textbf{Architecture implementation constraints.}
          We use the \texttt{fla-hub} port of RWKV-7 for HF/PEFT
          compatibility; it does not yet fully replicate the
          native CUDA reference's initialisation rules and
          layer-specific learning rates, so performance may sit
          slightly below the theoretical maximum.
    \item \textbf{Pre-training bias.} Both systems inherit biases
          from web-scale pre-training corpora; predictions may
          reflect societal stereotypes present upstream.
    \item \textbf{Closed-task caveat.} No external data used in
          fine-tuning, but pre-training corpora are \emph{not}
          controlled or audited.
    \item \textbf{Label schema.} Binary entailment is a
          simplification — real-world inference often includes a
          \emph{neutral} class.
    \item \textbf{Intended use.} Research and coursework only;
          not validated for high-stakes domains (legal, medical,
          safety-critical).
\end{itemize}
\end{block}

\vfill
\end{minipage}%

%==============================================================
%  FOOTER
%==============================================================
\vspace{0.3cm}
\begin{center}
\textcolor{uomPurple}{\rule{0.95\paperwidth}{1pt}}\\[0.2cm]
{\small
\faGithub\ \ \texttt{https://github.com/d4k3r/NLU} \quad
\faFile\ \ Model cards \& demo notebooks in Canvas submission
}
\end{center}

\end{frame}
\end{document}
